\section{Courbes elliptiques complexes}










\subsection{Courbes elliptiques complexes et tores}

Dans cette section, toutes les courbes seront définie sur $\Kb = \C$. Dans ce contexte, On démontre l'existence d'un point de vue analytique des courbes elliptiques :

\begin{Def}
	Un réseau $\Lambda$ est un ensemble $\left\lbrace nz_1 + mz_2 : n,m\in\Z\right\rbrace $ où $z_1 ,z_2$ sont deux complexes linéairement indépendant sur $\R$. Un tore complexe est un quotient $\C/\Lambda$ où $\Lambda$ est un réseau. 
	
\end{Def}

On montre comment associer à tout courbe elliptique définie sur $\C$ un réseau, et réciproquement. Il y a alors une correspondance entre les isogénies et les morphismes de tores. 

\begin{Def}
	Soit $\Lambda$ un réseau. on lui associe la fonction de Weierstrass $\wp$ relative à $\Lambda$, définie par
	\begin{equation*}
		\wp(z,\Lambda) = \frac{1}{z^2} + \sum_{\omega\in\Lambda\setminus\left\lbrace 0\right\rbrace }\left( \frac{1}{\left( z - \omega^2\right)} - \frac{1}{\omega^{2}} \right) 
	\end{equation*}
	On lui associe aussi pour tout entier $k>1$ sa série d'Eisenstein de poids $2k$ défini par
	\begin{equation*}
		G_{2k}(\Lambda) = \sum_{\omega\in\Lambda\setminus\left\lbrace 0\right\rbrace }\omega^{-2k}
	\end{equation*}
\end{Def}

On énonce les propriétés de ces objets qui nous serviront à énoncer le lien entre courbes elliptiques et tores.

\begin{Prop}\cite{Sil09}[VI 3.1, 3.5]
	
	Soit $\Lambda$ un réseau.
	\begin{enumerate}
		\item Pour tout entier $k>1$, la série $G_{2k}(\Lambda)$ est absolument convergente
		\item $\wp(z,\Lambda)$ vérifie l'équation différentielle
		\begin{equation*}
			(\wp')^{2} = 4\wp^{3} - 60G_{4}\wp - 140G_{6}
		\end{equation*}
	\end{enumerate}
\end{Prop}

\begin{Def}
	Soit $\Lambda$ un réseau. On utilise les notations standard $g_{2}(\Lambda)$ et $g_{3}(\Lambda)$ pour désigner les valeurs respectives de $G_{4}(\Lambda)$ et $G_{6}(\Lambda)$.
\end{Def}

\begin{Prop}\cite{Sil09}[VI 3.6]
	
	Soit $\Lambda$ un réseau. L'équation $E : y^2 = 4x^2 - g_{2}x - g_{3}$ défini un modèle de courbe elliptique (qui n'est pas sous forme de Weierstrass). L'application
	\begin{equation*}
		\psi : \C/\Lambda \rightarrow E(\C)\; ,\; z\mapsto \left(\wp(z), \wp'(z)\right)
	\end{equation*}
	est un isomorphisme. 
\end{Prop}

\begin{Rem}
	L'isomorphisme énoncé est au sens de groupe de Lie complexe. Nous utiliseront en particulier le fait que c'est un isomorphisme de groupe. 
\end{Rem}

\begin{Def}
	Soit $\Lambda$ un réseau. On note $E_{\Lambda}$ la courbe elliptique dont d'équation $E : y^{2} = 4x^{2} - g_2 x - g_3$ est un modèle. 
\end{Def}

Le théorème suivant, appelé théorème d'uniformisation, permet réciproquement d'associer à toute courbe elliptique complexe un réseau défini à homothétie prêt.

\begin{The}{Théorème d'uniformisation}\cite{Sil09}[VI, 5.1]
	
	Soit $A$, $B\in\C$ tels que $4a^3 - 27b^2\neq0$. Il existe un unique réseau $\Lambda$ tel que $g_{2}(\Lambda) = A$ et $g_{3}(\Lambda) = B$. 
\end{The}

\begin{Coro}\cite{Sil09}[VI, 5.2]
	
	Soit $E$ une courbe elliptique définie sur $\C$. Il existe un réseau $\Lambda$ unique à homothétie prêt tel que l'application
	\begin{equation*}
		\psi : \C/\Lambda \rightarrow E(\C)\; ,\; z\mapsto \left(\wp(z), \wp'(z)\right)
	\end{equation*}
	soit un isomorphisme entre le tore $\C/\Lambda$ et un modèle de $E$.
\end{Coro}






\subsection{Endomorphisme de tores et normalisation d'endomorphisme}

\begin{Def}
	Soit $\C/\Lambda_1$, $\C/\Lambda_2$ deux tores, et $\alpha\in\C$ tel que
	\begin{equation*}
		\alpha\in\left\lbrace \alpha\in\C : \alpha\Lambda_1 \subset \Lambda_2\right\rbrace  
	\end{equation*} 
	Un morphisme de tores est une application de la forme
	\begin{equation*}
		\phi_{\alpha} : \C/\Lambda_1 \rightarrow \C/\Lambda_2\; ,\; z \mapsto \alpha z
	\end{equation*}
	
\end{Def}

\begin{Rem}
	Si $\phi: \C/\Lambda_1 \rightarrow \C/\Lambda_2$ est une application holomorphe et nulle en $0$, alors elle de la forme $\phi_{\alpha}$ \cite{Sil09}[VI,4.1] pour un certain $\alpha$. On choisit donc de définir les morphismes de tores comme étant les applications de la forme $\phi_{\alpha}$.
	
\end{Rem}

L'isomorphisme entre un réseau $\Lambda$ et sa courbe associée $E(\Lambda)$ induit un isomorphisme entre anneaux d'endomorphismes.

\begin{The}\cite{Sil09}[VI 4.1, 5.3]
	
	Les catégories suivantes sont équivalentes
	\begin{enumerate}
		\item Les courbes elliptiques complexes, munies de leurs isogénies
		\item Les réseaux à homothéties prêt, munis des morphismes de tores
	\end{enumerate}
\end{The}

\begin{Coro}
	Soit $\Lambda$ un réseau. Soit $R\in\C$ l'anneau des complexes représentant les endomorphisme du tore $\C/\Lambda$
	\begin{equation*}
		R = \left\lbrace \phi_{\alpha} :  \alpha\in\C ,\; \alpha\Lambda \subset \Lambda \right\rbrace 
	\end{equation*}
	
	Alors on a l'isomorphisme d'anneau $R\simeq End(E_{\Lambda})$
\end{Coro}

\begin{Coro}\cite{Sil09}[VI, 4.1.1]
	Soit $E$, $E'$ deux courbes elliptiques. $E = E'$ si et seulement si leurs réseaux associés sont homothétiques.
\end{Coro}


Lorsque $E$ est une courbe elliptique à multiplications complexes, $End(E)$ peut être inclus dans $\C$ de deux manières (à conjugaison prêt dans le corps quadratique imaginaire). On souhaite normaliser le choix de l'injection de $End(E)$ dans $\C$. Pour cela, on veut associer à tout $\alpha$ tel que $\alpha\Lambda\subset\Lambda$, un endomorphisme notée $[\alpha]$. Le corollaire 3.6 fourni une première convention. On va cependant définir une seconde convention de normalisation à l'aide de l'invariant différentiel, qui se généralisera plus facilement pour les courbes elliptiques à multiplication complexe définie sur des corps finis. On montrera que les deux conventions coïncident.

\begin{Def}
	Soit $E : y^2 = x^3 + ax + b$ un modèle d'une courbe elliptique $E$, et $\phi : E\rightarrow E$. $\phi$ admet une unique forme explicite dont le départ et l'arrivé sont le modèle $E : y^2 = x^3 + ax + b$. Soit $\omega$ l'invariant différentiel de ce modèle. Il existe une constante $c\in\C$ telle que l'application contravariante de cette forme explicite de $\phi$ vérifie 
	\begin{equation*}
		(\phi)^{*}(\omega) = c\omega
	\end{equation*}
	La constante $c$ ne dépend pas du modèle de $E$ choisi, et est appelée le facteur de l'endomorphisme $\phi$.
\end{Def}

On souhaite construire un morphisme d'anneau $\alpha\mapsto\left[ \alpha\right]$ de sorte que le facteur de $\left[ \alpha\right]$ soit $\alpha$. On applique pour cela le corollaire 2.16 au cas où $\K = \C$.

\begin{Prop}\cite{Sil09}[III, 5.6]
	
	Soit $E$ une courbe elliptique sur $\C$, $\omega$ un invariant différentiel d'un modèle de $E$. Pour tout endomorphisme $\phi\in End(E)$, on note $a_{\phi}\in\Kb$ tel que $\phi^{*}(\omega) = a_{\phi}\omega$. Alors le morphisme d'anneau
	
	\begin{equation*}
		f : End(E)\rightarrow\Kb , \phi\longmapsto a_{\phi}
	\end{equation*}
	est injectif.
	
\end{Prop}


Soit $\OR$ l'image de $f$ dans la proposition précédente. L'application $[.] : \OR \rightarrow End(E)$ recherchée la réciproque de $f$. 

\begin{Prop}
	Soit $E$ une courbe elliptique sur $\C$ à multiplication complexe. Il existe un unique isomorphisme d'anneaux
	\begin{equation*}
		[.] : \OR\subset\C\rightarrow End(E)
	\end{equation*}
	tel que pour tout invariant différentiel $\omega$ d'un modèle de $E$,
	\begin{equation*}
		\forall\alpha\in\OR\;,\; [\alpha]^{*}(\omega) = \alpha\omega
	\end{equation*}
\end{Prop}

\begin{proof}
	On choisit $\OR$ comme l'image de l'application $f$ de la proposition 3.8. Alors la réciproque de $f$ convient. Le facteur de tout endomorphisme étant unique, un tel isomorphisme est unique. 
\end{proof}

\begin{Def}
	Soit $E$ une courbe elliptique définie sur $\C$ à multiplication complexe. Soit $[.] : \OR\subset\C\rightarrow End(E)$ un isomorphisme d'anneaux. Le couple $(E,[.])$ est normalisé si pour tout invariant différentiel $\omega$ d'un modèle de $E$.
	\begin{equation*}
		\forall\alpha\in\OR\;,\; [\alpha]^{*}(\omega) = \alpha\omega
	\end{equation*}
\end{Def}

On montre que avec cette convention de normalisation, l'image $\OR$ de $End(E)$ dans $\C$ correspond à l'anneau d'endomorphisme de tous tores équivalent à $E$ :


\begin{Prop}\cite{Sil94}[II 1.1]
	
	Soit $(E,[.])$ une courbe elliptique définie sur $\C$ à multiplication complexe normalisé. Soit $[\alpha]$ un endomorphisme de facteur $\alpha$. Si $\Lambda$ est un réseau associé à $E$, alors $\alpha\Lambda\in\Lambda$ et l'endomorphisme $[\alpha]$ est équivalent à l'endomorphisme de tore $\phi_{\alpha}$.
\end{Prop}

La donnée d'un réseau associé à une courbe elliptique $E$ à multiplication complexe permet d'expliciter le corps quadratiques imaginaire dont l'anneau d'endomorphisme est un ordre.

\begin{Prop}
	Soit $\Lambda = \left[ z_{1}, z_{2}\right] $ un réseau , et $E_{\Lambda}$ sa courbe elliptique associée. Si $E_{\Lambda}$ est à multiplication complexe, alors $\K = \Q\left( \frac{z_1}{z_2}\right)$ est un corps quadratique imaginaire dont $End(E)$ est un ordre. 
\end{Prop}

\begin{proof}\cite{Sil09}[VI 5.5]
	
	Soit $\tau = \frac{z_1}{z_2}$. $\Lambda$ est homothétique au réseau engendré par 1 et $\tau$, donc on peut supposer que $z_1 = 1$ et $z_2 = \tau$ sans perte de généralité. Soit $R =  \left\lbrace\alpha\in\C, \alpha\Lambda\subset\lambda\right\rbrace$, de sorte que $End(E)\simeq R$. Tout $\alpha\in R$ vérifie $\alpha\in\Lambda$ et $\alpha\tau\in\Lambda$. On en déduit que $\alpha$ est racine d'un polynôme de degrés deux à coefficient dans $\Z$, donc les éléments de $R$ sont des entiers algébriques. De plus il existe $\alpha\in R$ qui n'est pas dans $\Z$. On en déduit l'existence d'un polynôme annulateur de degrés 2 pour $\tau$, qui n'est pas dans $\R$. Donc $\Q\left(\tau\right)$ est un corps quadratique imaginaire dont $R$ est un ordre. 
\end{proof}








\subsection{Courbes d'anneaux d'endomorphismes donné}

Dans cette section, les preuves seront rédigés. Soit $E_{\Lambda}$ une courbe elliptique associée à un réseau $\Lambda$.  On supposera toujours l'isomorphisme $End(E)\simeq\OR_{\Delta}$ normalisé. Dans la suite, $\K\subset\C$ est le corps quadratiques imaginaire dont $\OR_{\Delta}$ est un ordre.

\begin{Def}
	Soit $\OR_{\Delta}$ un ordre de discriminant $\Delta$ d'un corps quadratique imaginaire. On note $Ell_{\C}(\OR_{\Delta})$ l'ensemble des courbes elliptiques à multiplication complexes d'anneaux d'endomorphismes $\OR_{\Delta}$. 
\end{Def}

Partant d'un ordre $\OR_{\Delta}$, on peut trouver des réseaux $\Lambda$ tels que $End(E_{\Lambda}) = \OR_{\Delta}$. En effet, soit $\LR$ un idéal fractionnaire propre de $\OR_{\Delta}$. Un tel idéal est un $\Z$-module libre de rang $2$ qui n'est pas inclus dans $\R$, donc un réseau. On peut donc lui associer une courbe elliptique $E_{\LR}$.

\begin{Prop}
	Soit $\OR_{\Delta}$ un ordre de discriminant $\Delta$ d'un corps quadratique imaginaire. Soit $\LR\subset\OR_{\Delta}$ un idéal fractionnaire propre. Alors $End(E_{\LR})\simeq\OR_{\Delta}$, donc $E_{\LR}\in Ell_{\C}(\OR_{\Delta})$.
\end{Prop}

\begin{proof}\cite{Sil94}[II, 1]
	
	Par hypothèse, $\LR$ est fractionnaire propre, donc :
	\begin{align*}
		End(E_{\LR})&\simeq\left\lbrace\alpha\in\C : \alpha\LR\subset\LR\right\rbrace\\
		&= \left\lbrace\alpha\in\K : \alpha\LR\subset\LR\right\rbrace = \OR_{\Delta}
	\end{align*}
\end{proof}

Puisque des réseaux homothétiques définissent des courbes équivalentes, pour tout $c\in\K$, $c\LR$ est à la fois un idéal de $\OR_{\Delta}$ et un réseau homothétique à $\LR$. On s'intéresse donc aux classes d'idéaux. Soit $Cl(\OR_{\Delta})$ le groupe de classes de $\OR_{\Delta}$, et $\CL$ la classe de l'idéal fractionnaire propre $\LR$. L'application
\begin{equation*}
	Cl(\OR_{\Delta})\rightarrow Ell_{\C}(\OR_{\Delta}),\; \CL\mapsto E_{\LR}
\end{equation*} 
est donc bien définie. On va montrer qu'il s'agit d'une bijection

\begin{Def}
	Soit $\Lambda$ un réseau et $\LR$ un idéal fractionnaire propre de $\OR_{\Delta}$. On défini le produit $\LR\Lambda$ par
	\begin{equation*}
		\LR\Lambda = \left\lbrace \alpha_1\lambda_1 + \ldots + \alpha_r\lambda_r : \alpha_i\in\LR , \lambda_i\in\Lambda , r\in\N\right\rbrace 
	\end{equation*}
\end{Def}

\begin{Prop}
	Soit $\Lambda$ un réseau et $\LR$ un idéal fractionnaire propre de $\OR_{\Delta}$. Alors $\LR\Lambda$ est un réseau de $\C$. De plus, si $E_{\Lambda}\in Ell_{\C}(\OR_{\Delta})$ alors $E_{\LR\Lambda}\in Ell_{\C}(\OR_{\Delta})$
\end{Prop}

\begin{proof}\cite{Sil94}[II 1.2]
	
	Puisque $End(E_{\Lambda}) \simeq \OR_{\Delta}$ est normalisé,  $\OR_{\Delta}\Lambda = \Lambda$. Il existe $d\in\Z$ tel que $d\LR\subset\OR_{\Delta}$. Alors $\LR\Lambda\subset\frac{1}{d}\Lambda$ donc $\LR\Lambda$ est inclus dans un réseau de $\C$. 
	Maintenant soit $a\in\Z\cap\LR$. Alors $a\OR_{\Delta}\subset\LR$ donc $a\Lambda\subset\LR\Lambda$ car $\OR_{\Delta}\Lambda = \Lambda$. Ainsi $\LR\Lambda$ contient un réseau de $\C$. $\LR\Lambda$ est compris entre deux réseaux, donc est lui même un réseau de $\C$.
	Pour montrer que $E_{\LR\Lambda}\in Ell_{\C}(\OR_{\Delta})$ il suffit de montrer que $End(E_{\LR\Lambda}) = End(E_{\Lambda})$, ce qui est clair car $\LR$ est inversible. 
\end{proof}

Soit $\Lambda$ un réseau tel que $End(E_{\Lambda}) \simeq \OR_{\Delta}$. Pour tout idéal fractionnaire propre, $End(E_{\LR\Lambda}) = End(E_{\Lambda})$, mais $\Lambda$ et $\LR\Lambda$ ne sont pas à priori homothétique. $E_{\Lambda}$ et $E_{\LR\Lambda}$ sont donc deux éléments distincts de $Ell_{\C}(\OR_{\Delta})$. 

\begin{Prop}
	Soit $\Lambda$ un réseau tel que $E_{\Lambda}\in Ell_{\C}(\OR_{\Delta})$, et $\LR$, $\LR'$ deux idéaux fractionnaires propres de $\OR_{\Delta}$. 
	\begin{enumerate}
		\item $E_{\LR\Lambda} = E_{\LR'\Lambda}$ si et seulement si $\CL = \CL'$
		\item Le groupe $C(\OR_{\Delta})$ agit sur $Ell_{\C}(\OR_{\Delta})$ par $\CL\cdot E_{\Lambda} = E_{\LR^{-1}\Lambda}$
		\item Cette action est simplement transitive. En particulier $\left| C(\OR_{\Delta})\right| = \left| Ell_{\C}(\OR_{\Delta})\right| $
	\end{enumerate}
\end{Prop}

\begin{proof}\cite{Sil94}[II,1.2]
	
	\begin{enumerate}
		\item On sait que la normalisation de $End(E)$ implique que $\OR_{\Delta}$ est l'ensemble des facteurs des endomorphismes de $E$, donc $\OR_{\Delta}\Lambda = \Lambda$ d'après la proposition 3.10. De plus $E_{\LR\Lambda} \simeq E_{\LR'\Lambda}$ si et seulement s'il existe $c\in\C^{*}$ tel que $\LR\Lambda = c\LR'\Lambda$. $\LR$ et $\LR'$ étant inversible, cela équivaut à 
		\begin{equation*}
			\Lambda = c\LR^{-1}\LR'\Lambda = c^{-1}\LR'^{-1}\LR\Lambda
		\end{equation*}
		Donc $c\LR^{-1}\LR'$ et $c^{-1}\LR'^{-1}\LR$ fixent $\Lambda$, donc ils sont inclus dans $\OR_{\Delta}$. Étant inverse l'un de l'autre, le seul idéal entier de $\OR_{\Delta}$ dont l'inverse est aussi entier est $\OR_{\Delta}$ lui même. Donc $c\LR^{-1}\LR' = \OR_{\Delta}$ et $\LR = c\LR'$
		
		\item L'action est bien définie sans difficulté, la présence de l'inverse dans la définition n'est pas nécessaire pour définir l'action, mais sera utile par la suite. 
		
		\item On se donne deux réseaux $\Lambda_1$ et $\Lambda_2$ dont les courbes elliptiques associées sont dans $Ell_{\C}(\OR_{\Delta})$, et on cherche $\CL$ tel que $\CL\cdot E_{\Lambda_1} = E_{\Lambda_2}$. Cela équivaut à trouver un idéal fractionnaire $\LR$ tel que $\Lambda_2$ et $\LR^{-1}\Lambda_1$ soient homothétiques. Mais puisque $E_{\Lambda_1}\in Ell_{\C}(\OR_{\Delta})$, $\OR_{\Delta}\Lambda_1 = \Lambda_1$ et 
		\begin{equation*}
			\left\lbrace \alpha\in\C : \alpha\Lambda_1\subset\Lambda_1\right\rbrace = \OR_{\Delta}
		\end{equation*}
		De plus 
		\begin{equation*}
			\forall \beta\in\C^{*}\;,\; \alpha\Lambda_1\subset\Lambda_1\iff\alpha\beta\Lambda_1\subset\beta\Lambda_1
		\end{equation*}
		Donc si l'on trouve $\beta$ tel que $\beta\Lambda_1\subset\K$, alors $\beta\Lambda_1$ sera un idéal fractionnaire propre de $\OR_{\Delta}$. Soit $z_1$, $z_2$ des générateurs de $\Lambda_1$. D'après la proposition 3.11, si $\beta = \frac{z_1}{z_2}$ alors $\beta\Lambda_1\subset\K$ où $\K$ est le corps quadratique imaginaire dont $\OR_{\Delta}$ est un ordre. Et $\OR_{\Delta}\beta\Lambda_1 = \beta\Lambda_1$ donc $\beta\Lambda_1$ est un $\OR_{\Delta}$-module de type fini, donc c'est un idéal fractionnaire propre de $\OR_{\Delta}$. Soit $\LR_1 = \beta\Lambda_1$, et $\LR_2$ définie de la même manière à partir de $\Lambda_2$. Alors $\LR = \LR_2\LR_1^{-1}$ convient : $\LR^{-1}\Lambda_1$ est homothétique à $\Lambda_2$.
		
		\item Pour finir, la simple transitivité découle du point 1.
		
	\end{enumerate}
	
\end{proof}

Notons que l'on peut normaliser toutes les courbes elliptiques de $Ell_{\C}(\OR_{\Delta})$ de façon cohérente.

\begin{Prop}
	Soit $(E_1,[.]_1)$ et $(E_2,[.]_2)$ deux courbes elliptiques définies sur $\C$, à multiplication complexe, normalisées et de même anneaux d'endomorphisme $\OR_{\Delta}$. Soit $\phi : E_1 \rightarrow E_2$ une isogénie. Alors 
	\begin{equation*}
		\forall\alpha\in\OR_{\Delta},\;\phi\circ[\alpha]_1 = [\alpha]_2\circ\phi
	\end{equation*}
\end{Prop}

\begin{proof}\cite{Sil94}[II, 1.1.1]
	
	Si $\omega$ est un invariant différentiel de $E_2$, alors $\phi^{*}(\omega)$ en est un pour $E_1$. $[\alpha]_{1}^{*}$ agit dessus par multiplication par $\alpha$. Ainsi
	\begin{align*}
		(\phi\circ[\alpha]_1)^{*}(\omega) &= [\alpha]_{1}^{*}(\phi^{*}(\omega))\\
		&= \alpha\phi^{*}(\omega)\\
		&= \phi^{*}(\alpha\omega)\\
		&= \phi^{*}([\alpha]_{2}^{*}(\omega))\\
		&= ([\alpha]_2\circ\phi)^{*}(\omega)\\
	\end{align*}
	Donc 
	\begin{equation*}
		(\phi\circ[\alpha]_1)^{*} = ([\alpha]_2\circ\phi)^{*}
	\end{equation*}
	D'après la proposition 2.8, puisque toutes les isogénies sont séparable en caractéristique 0, l'application qui à une isogénie associe son application contravariante est injective. Donc 
	\begin{equation*}
		\phi\circ[\alpha]_1 = [\alpha]_2\circ\phi
	\end{equation*}
\end{proof}

Si $\phi : E \rightarrow E'$ est une isogénie entre deux courbes elliptiques $E$, $E'$ distinctes de même anneaux d'endomorphisme $\OR_{\Delta}$, alors $\phi$ correspond à l'action d'un unique élément $\CL\in C(\OR_{\Delta})$ sur $E$, puisque l'action est simplement transitive. On ne pourra pas déterminer $\phi$ à partir de $\CL$, car $\CL$ défini $\phi$ à endomorphisme prêt au départ et à l'arrivé (c'est dire à action d'un idéal fractionnaire prêt). 










\subsection{Groupes de torsions d'idéaux}

On se donne une courbe elliptique $E\in Ell_{\C}(\OR_{\Delta})$ et un idéal fractionnaire propre $\LR$. On souhaite définir à partir de $\LR$ une isogénie $\phi : E\rightarrow \CL\cdot E$. Cela équivaut à définir le noyau de $\phi$ à partir de $\LR$. Lorsque $\LR$ est un idéal entier propre, on va définir un sous groupe fini de $E$ associé à $\LR$.

\begin{Def}
	Soit $(E,[.])$ une courbe elliptique définie sur $\C$ à multiplication complexe par $\OR_{\Delta}$, normalisée. Soit $\LR\in\OR_{\Delta}$ un idéal entier propre. Le groupe de $\LR$-torsion d'un modèle de $E$ est 
	\begin{equation*}
		E\left[ \LR\right] = \left\lbrace P\in E : \left[ \alpha\right] (P) = \OR \; \forall\alpha\in\LR \right\rbrace 
	\end{equation*}
	Le groupe de $\LR$-torsion de $E$ est définie à isomorphisme prêt.
\end{Def}


%ICI


\begin{Ex}
	\begin{itemize}
		
		\item Si $\LR = m\OR_{\Delta}$ avec $m\in\Z$, alors $E\left[ \LR\right] = E\left[ m\right] $ est le groupe de $m$-torsion.
		
		\item Si $\alpha\in\OR_{\Delta}$, $E[\alpha\OR_{\Delta}] = Ker([\alpha])$ (se démontre par double inclusions).
		
	\end{itemize}
	
\end{Ex}

On montre que $E[\LR]$ est le noyau de l'isogénie $\phi : E\rightarrow \CL\cdot E$. La présence d'un inverse dans la définition de $\CL\cdot E$ permet ce résultat.

\begin{Def}
	Soit $\LR$ un idéal entier propre de $\OR_{\Delta}$, et $\Lambda$ un réseau tel que $E_{\Lambda}\in Ell_{\C}(\OR_{\Delta})$. On associe à $\LR$ le morphisme de tores
	\begin{equation*}
		\phi_{\LR} : \C/\Lambda\rightarrow\C/\LR^{-1}\Lambda\; , \; z\mapsto z
	\end{equation*}
	
	
\end{Def}

\begin{Prop}
	Soit $\Lambda$ un réseau tel que $E_{\Lambda}\in Ell_{\C}(\OR_{\Delta})$, et $\LR\in\OR_{\Delta}$ un idéal entier propre. Alors $E_{\Lambda}\left[ \LR\right]$ est un sous groupe fini de $E_{\Lambda}$, et $E_{\LR^{-1}\Lambda} = \CL\cdot E_{\Lambda}$ est l'unique courbes elliptiques telle qu'il existe une isogénie $\phi : E_{\Lambda}\rightarrow E_{\LR^{-1}\Lambda}$ de noyau $E_{\Lambda}\left[ \LR\right]$. Cette isogénie est $\phi_{\LR}$
	
\end{Prop}


\begin{proof}\cite{Sil94}[II,1.4]
	
	On montre que $E_{\Lambda}\left[ \LR\right]$ est le noyau de $\phi_{\LR}$. D'abord, en identifiant un endomorphisme à son facteur,
	\begin{align*}
		E_{\Lambda}\left[ \LR\right]&\simeq \left\lbrace z\in\C/\Lambda : \alpha z = 0 \;,\; \forall\alpha\in\LR\right\rbrace \\
		&=\left\lbrace z\in\C : \alpha z\in\Lambda \;,\;\forall\alpha\in\LR\right\rbrace / \Lambda\\
		&=\left\lbrace z\in\C : z\LR\subset\Lambda\right\rbrace / \Lambda
	\end{align*}
	Or $\LR\Lambda = \Lambda$ car $\LR$ est entier, donc $ \LR^{-1}\Lambda \subset \left\lbrace z\in\C : z\LR\subset\Lambda\right\rbrace $. Inversement si $z\in\C$ est tel que $z\LR\subset\Lambda$, alors $z\OR_{\Delta}\subset\LR^{-1}\Lambda$ donc en particulier $z\in\LR^{-1}\Lambda$. Donc $ \LR^{-1}\Lambda \subset \left\lbrace z\in\C : z\LR\subset\Lambda\right\rbrace $ et
	\begin{equation*}
		E_{\Lambda}\left[ \LR\right]\simeq\LR^{-1}\Lambda/\Lambda = Ker(\phi_{\LR})
	\end{equation*}
	En particulier $E_{\Lambda}\left[ \LR\right]$ est un sous groupe fini de $E$. 
	
\end{proof}

Lorsque $\LR\subset\OR_{\Delta}$ n'est plus seulement un idéal entier propre mais est un idéal premier avec le conducteurs, le degré de l'isogénie $\phi_{\LR}$ est la norme de $\LR$. 

\begin{Prop}
	Soit $\Lambda$ un réseau tel que$E_{\Lambda}\in Ell_{\C}(\OR_{\Delta})$, et $\LR\in\OR_{\Delta}$ un idéal entier premier avec le conducteur. Alors :
	\begin{enumerate}
		\item $E\left[\LR\right]$ est un $\OR_{\Delta}/\LR$-module libre de rang $1$.
		\item $\phi_{\LR}$ est de degré $N\left( \LR\right)$ 
		\item Si $\alpha\in\OR_{\Delta}$, $[\alpha]$ est de degré $N_{\Q}^{\K}\left( \alpha\right) $
	\end{enumerate}
	
\end{Prop}

\begin{proof}\cite{Sil94}[II, 1.4, 1.5]
	
	\begin{enumerate}
		\item (Admis) La preuve de \cite{Sil94} utilise une version du théorème des restes chinois sur les idéaux fractionnaire, sur l'anneau des entiers $\OR_{\K}$. La preuve ne se généralise donc aux ordres quelconques $\OR_{\Delta}$ qu'en se restreignant aux idéaux entier premier avec le conducteur.  
		
		\item Découle de 1. car $\left| E\left[\LR\right]\right| = \left| \OR_{\Delta}/\LR\right| = N\left( \LR\right)$.
		
		\item En particulier si $\LR = \alpha\OR_{\Delta}$ est principal, $[\alpha]$ est l'endomorphisme de noyau $E\left[\alpha\OR_{\Delta}\right]$. Donc 
		\begin{equation*}
			deg \left( \left[\alpha \right] \right) = N\left( \alpha\OR_{\Delta}\right) = N_{\Q}^{\K}\left( \alpha\right) 
		\end{equation*}
	\end{enumerate}
\end{proof}


Ainsi, se donner un idéal entier propre $\LR$ de $\OR_{\Delta}$ et une courbe elliptique $E$ d'anneaux d'endomorphisme $\OR_{\Delta}$ permet de définir une isogénie $\phi_{\LR}$ et de trouver la classe d'équivalence des courbes d'arrivée. Il suffit alors d'appliquer un isomorphisme à l'arrivée pour que l'isogénie soit normalisé au sens de la proposition 2.13. 



\begin{Rem}
	Si $E$ est une courbe elliptique de $Ell_{\C}(\OR_{\Delta})$, alors tout idéal fractionnaire propre de $\OR_{\Delta}$ défini une isogénie. On est seulement capable d'expliciter le noyau quand l'idéal est entier. En particulier tout idéal fractionnaire principal défini un endomorphisme de $E$, lui même toujours représenter par un idéal principal entier défini par son facteur. Donc deux idéaux factionnaires propres peuvent définir une même isogénie. 
\end{Rem}